\documentclass{article}
\usepackage{graphicx}
\usepackage{amsmath}
\usepackage{amssymb}   

\begin{document}

\section*{Notes on sampling}
What we want to be able to do, is to obtain an effects matrix $B$ of the confounders on variables $X$ and $Y$. Now these effect
sizes should be chosen in a very particular way, to reflect how strongly the confounders affect the variables as a whole, and how
much of that effect is caused by a non-linear component of the confounders. 

\subsection*{Defining the parameters}
We take a number of base confounders, which exert a linear effect on the variables $X$ and $Y$. We call this number $k_{\text{lin}}$. Now every confounder
also has a number of non-linear components, which exert a non-linear effect on the variables $X$ and $Y$. We call the sum over the amount of these
variables for all confounders $k_{\text{nonlin}}$. The total number of confounders is therefore $k = k_{\text{lin}} + k_{\text{nonlin}}$. In the derivations below we will
first work this out for the case where each confounder has only one non-linear component. Later on we will extend this to multiple non-linear components.

\noindent
Now the $k$ confounders exert an effect on the variables $X$ and $Y$. These effects can be captured in the vectors:

\[
\gamma_{x} \in \mathbb{R}^k
\]

and

\[
\gamma_{y} \in \mathbb{R}^k
\]

\noindent
meaning the entries to both vectors can take on any real number. Both vectors together can be written as

\[
\theta = (\gamma_{x}, \gamma_{y}) \in \mathbb{R}^{2k}.
\]

\noindent
Now we constrain each coefficient to be between $0.01$ and $0.30$, since larger values seem unrealisitic and may lead to difficulties to get the system stationary.
However, we could always opt for a larger upper bound if needed. We can summarise this with the following expression:

\[
\Theta = \{\, \theta = (\gamma_x, \gamma_y) : \gamma_x, \gamma_y \in [0.01, 0.30]^k \,\}.
\]

\subsection*{The influence of confounders}

We can compute the total confounder induced variance in $X$ and $Y$ as follows: 

\[
S_U =
\begin{pmatrix}
\operatorname{Var}_X(U) & \operatorname{Cov}_{XY}(U) \\
\operatorname{Cov}_{XY}(U) & \operatorname{Var}_Y(U)
\end{pmatrix}
=
\begin{pmatrix}
\sum_j \gamma_{xj}^2 & \sum_j \gamma_{xj}\gamma_{yj} \\
\sum_j \gamma_{xj}\gamma_{yj} & \sum_j \gamma_{yj}^2
\end{pmatrix}.
\]

\noindent
Now when the confounders are standard normal and uncorrelated, we can simplify this to:

\[
\operatorname{Var}_X^U = \sum_{i=1}^{k} \gamma_{x,i}^2.
\]

\noindent
Adding any covariances just complicates this calculation a little bit, but is not all that difficult. Now we can write the total influence of the 
confounders as:

\[
S(\theta) = \sum_{i=1}^{k} \gamma_{x,i}^2 + \sum_{i=1}^{k} \gamma_{y,i}^2 = s_0.
\]

\noindent
Where later on, we want to pick $s_0$. 

\subsection*{Linear vs non-linear influence}
Now we want to compute which part of this variance is attributable to a linear component, and which part is attributable to a non-linear component.
This is a measure of non-linearity which falls between $0$ and $1$. We can partition our vectors of coefficients as follows:

\[
\gamma_x = (\gamma_{x,\text{lin}}, \gamma_{x,\text{nl}}) \qquad\qquad 
\gamma_y = (\gamma_{y,\text{lin}}, \gamma_{y,\text{nl}})
\]

\noindent
where the first $k_{\text{lin}}$ entries are the linear parts, and the last $k_{\text{nonlin}}$ entries are the non-linear parts. Now we can compute for both
$X$ and $Y$ the non-linear part of the explained variance as:

\[
\eta_X(\theta)= \sum_{i=k_{\text{lin}}+1}^{k} \gamma_{x,i}^2 \qquad\qquad 
\eta_Y(\theta) = \sum_{i=k_{\text{lin}}+1}^{k} \gamma_{y,i}^2.
\]

\noindent
Then their average gives us the non-linear influence as:

\[
\eta(\theta) = \frac{\eta_X + \eta_Y}{2} = \eta_0.
\]

\noindent
Where later we want to pick $\eta_0$.

\subsection*{Sampling procedure}
Now we can choose $\eta_0$ and $s_0$. We can also choose a number of linear confounders $k_{\text{lin}}$ and a number of 
non-linear components $k_{\text{nonlin}}$. We can call the number of non-linear components per confounder the complexity of the nonlinearity. 
As touched upon before, for now we will assume this complexity is $1$, meaning each confounder has exactly one non-linear component. Now
there is a set of parameters $\theta$ which satisfy these constraints. We can call the set of solutions to these constraints:

\[
M = \{\, \theta \in \Theta : S(\theta) = s_0,\ \eta(\theta) = \eta_0 \,\}.
\]

\noindent
Now the problem is, we do not know $M$, and as such we cannot sample from it. We can however sample from the set of all possible values in
our $2k$-dimensional space $\Theta$. Now we do something called rejection sampling, where we sample from this space and check if the 
constraints are met. If not, we reject the sample and sample again. Now, you might already see that this process will be incredibly inefficient,
as picking a random sample from $\Theta$ will almost never satisfy the constraints. 

There are some things we can do to make this process more efficient.
First of all, we can make the constraints a little bit less strict, by allowing a small margin of error around the desired values. This will ``thicken''
the set $M$ a little bit, making it more likely to hit a valid sample. Now, we also need to make sure that we pick our samples in a smarter way. In
practice this means choosing one more constraint. We will choose to divide up $s_0$ evenly over $X$ and $Y$. This means we add the constraint:

\[
\operatorname{Var}_X^U = \operatorname{Var}_Y^U = \frac{s_0}{2}.
\]

\noindent
Note that we could always choose a different way to divide up $s_0$ over $X$ and $Y$. We could each time also sample this division from a uniform distribution
before picking any values from the space of the parameters. In this way we would generalise over different ways to divide up the confounder influence 
over $X$ and $Y$. Note that we face a similar choice when dividing up $\eta_0$ over $X$ and $Y$. For now we will keep it simple and divide both evenly:

\[
\eta_X = \eta_Y = \eta_0.
\]

\noindent
Now recall the equation for $\eta_X$:

\[
\eta_X = \frac{N_X}{L_X + N_X},
\]

\noindent
which we can now rewrite as:

\[
N_X = \eta_0 (L_X + N_X).
\]

\noindent
Now note that the total amount of variance $L_X + N_X$ is equal to $s_0/2$. Therefore we can rewrite this as:

\[
N_X = \eta_0 \frac{s_0}{2},
\]

and

\[
N_Y = \eta_0 \frac{s_0}{2},
\]

and

\[
L_X = (1 - \eta_0) \frac{s_0}{2},
\]

and

\[
L_Y = (1 - \eta_0) \frac{s_0}{2}.
\]

\noindent
Now for each of these four quantities, we can calculate them using the sum of squares:

\[
L_X = \sum_{i=1}^{k_{\text{lin}}} \gamma_{x,\text{lin},i}^2,
\]

\[
L_Y = \sum_{i=1}^{k_{\text{lin}}} \gamma_{y,\text{lin},i}^2,
\]

\[
N_X = \sum_{i=1}^{k_{\text{nonlin}}} \gamma_{x,\text{nl},i}^2,
\]

\[
N_Y = \sum_{i=1}^{k_{\text{nonlin}}} \gamma_{y,\text{nl},i}^2.
\]

\noindent
Now this is where we can do something smart. We want to draw a random vector from any symmetric distribution, such that the sum of squares of the entries:

\[
u^{x,\text{lin}} \in \mathbb{R}^{k_{\text{lin}}},
\]

\noindent
we normalize this vector such that:

\[
\sum_{i=1}^{k_{\text{lin}}} (u^{x,\text{lin}}_i)^2 = 1,
\]

\noindent
and then we set:

\[
\gamma_{x,\text{lin},i} = \sqrt{L_X}\, u^{x,\text{lin}}_i,
\]

\noindent
which fulfills the constraint $L_X = \sum_{i=1}^{k_{\text{lin}}} \gamma_{x,\text{lin},i}^2$.  

We can do the same for the non-linear confounders:

\[
u^{x,\text{nl}} \in \mathbb{R}^{k_{\text{nonlin}}},
\]

\noindent
normalize:

\[
\sum_{i=1}^{k_{\text{nonlin}}} (u^{x,\text{nl}}_i)^2 = 1,
\]

\noindent
and set:

\[
\gamma_{x,\text{nl},i} = \sqrt{N_X}\, u^{x,\text{nl}}_i,
\]

\noindent
which fulfills the constraint $N_X = \sum_{i=1}^{k_{\text{nonlin}}} \gamma_{x,\text{nl},i}^2$.

We can do the same for $Y$. Now we also wanted our samples to fall
within the range between $0.01$ and $0.30$. We can simply check this after sampling and reject any samples that do not fall within this range.

To recap in words, we:

\begin{itemize}
    \item Choose $k_{\text{lin}}$, $k_{\text{nonlin}}$, $s_0$ and $\eta_0$
    \item Sample four random vectors from a symmetric distribution: $u^{x,\text{lin}}$, $u^{x,\text{nl}}$, $u^{y,\text{lin}}$, $u^{y,\text{nl}}$
    \item Normalize these vectors to have sum of squares equal to $1$
    \item Scale these vectors to have sum of squares equal to their respective budgets
    \item Combine these vectors to form $\gamma_x$ and $\gamma_y$
    \item Check if all entries of $\gamma_x$ and $\gamma_y$ fall within $[0.01, 0.30]$. If not, reject and start over.
\end{itemize}

\end{document}

% pdflatex Notes_on_sampling.tex 
% Remove-Item "*.aux","*.log","*.toc","*.out","*.synctex.gz" -Force